% RQ4 Methodology: Treasury Resilience Experiments

\subsection{Experimental Design}

We conduct a factorial experiment varying three treasury policy parameters:

\begin{table}[htbp]
\centering
\caption{RQ4 Experiment configuration (Experiment 06)}
\label{tab:rq4-experiment}
\begin{tabular}{@{}llll@{}}
\toprule
Factor & Levels & Values & Description \\
\midrule
Buffer fraction & 4 & 0\%, 10\%, 20\%, 30\% & Reserved portion \\
Spending limit & 3 & None, 5\%, 2\% of treasury & Per-period cap \\
Top-up mode & 3 & None, Revenue divert, Freeze & Emergency mechanism \\
\bottomrule
\end{tabular}
\end{table}

Total configurations: 12 (4 buffer $\times$ 3 spending limit, with top-up as modifier)

Each configuration runs 25 times with different seeds.

Total runs: 300

\subsection{Parameter Selection}

\subsubsection{Buffer Fraction Levels}

\begin{itemize}
    \item \textbf{0\%}: No reserve; all treasury available for spending
    \item \textbf{10\%}: Light buffer; protection against small shocks
    \item \textbf{20\%}: Moderate buffer; standard corporate practice parallel
    \item \textbf{30\%}: Conservative buffer; substantial protection but idle capital
\end{itemize}

\subsubsection{Spending Limit Levels}

\begin{itemize}
    \item \textbf{None}: Unlimited spending (subject to treasury availability)
    \item \textbf{5\% of treasury per period}: Moderate limit
    \item \textbf{2\% of treasury per period}: Conservative limit
\end{itemize}

\subsubsection{Top-Up Modes}

\begin{itemize}
    \item \textbf{None}: No emergency mechanism; buffer depletes without restoration
    \item \textbf{Revenue divert}: Redirect incoming revenue to restore buffer
    \item \textbf{Spending freeze}: Halt discretionary spending until buffer restored
\end{itemize}

\subsection{Market Conditions}

We simulate three market regimes:

\begin{table}[htbp]
\centering
\caption{Market regime parameters}
\label{tab:market-regimes}
\begin{tabular}{lcc}
\toprule
Regime & Drift ($\mu$) & Volatility ($\sigma$) \\
\midrule
Bull & +0.05\%/day & 2\%/day \\
Bear & -0.05\%/day & 4\%/day \\
Stable & 0\%/day & 1\%/day \\
\bottomrule
\end{tabular}
\end{table}

Regime changes occur stochastically (5\% probability per period).

\subsection{Baseline Configuration}

Common baseline:

\begin{itemize}
    \item Initial treasury: 10,000 units (normalized)
    \item Token concentration: 70\% native token, 30\% stablecoin
    \item Revenue rate: 0.5\% of treasury per period (average)
    \item Proposal spending demand: 1-3\% of treasury per period
    \item Simulation length: 2,000 steps
\end{itemize}

\subsection{Hypotheses}

\begin{itemize}
    \item \textbf{H4.1}: Higher buffer fractions reduce treasury volatility
    \item \textbf{H4.2}: Spending limits extend runway but reduce proposal throughput
    \item \textbf{H4.3}: Emergency top-ups reduce maximum drawdown severity
    \item \textbf{H4.4}: Buffer effectiveness depends on market regime (more valuable in bear markets)
\end{itemize}

\subsection{Metrics}

Primary metrics for RQ4:

\begin{description}
    \item[Final Treasury] Value at simulation end
    \item[Volatility] Standard deviation of treasury value
    \item[Max Drawdown] Largest peak-to-trough decline
    \item[Runway] Average periods until depletion (if applicable)
    \item[Top-up Frequency] Number of emergency top-up events
    \item[Proposal Throughput] Proposals funded per period
\end{description}

\subsection{Statistical Analysis}

\subsubsection{Regression Analysis}

Model treasury outcomes as function of policy parameters:

\begin{equation}
\text{Outcome} = \beta_0 + \beta_1 \cdot \text{buffer} + \beta_2 \cdot \text{limit} + \beta_3 \cdot \text{top-up} + \epsilon
\end{equation}

\subsubsection{Regime-Conditional Analysis}

Compare policy effectiveness across market regimes using interaction terms.

\subsubsection{Survival Analysis}

Model time-to-depletion as function of policy parameters using Cox proportional hazards.
